\chapter{Solución propuesta}\label{chap:2}

El presente capítulo expone la solución propuesta para el sistema FoT (\textit{Farm of Things}),
abordando la documentación de ingeniería de requisitos que sustenta su diseño. Se describen
el negocio, los requisitos funcionales, los atributos de calidad, las restricciones del sistema
y los casos de uso identificados como significativos para la arquitectura.

% ─────────────────────────────────────────────────────────────────────────────
\section{Documentación de Ingeniería de Requisitos}

La ingeniería de requisitos constituye la base sobre la cual se define y delimita el sistema.
A continuación se presenta la descripción del negocio, seguida de los requisitos funcionales,
los atributos de calidad y las restricciones que condicionan la solución.

% ─────────────────────────────────────────────────────────────────────────────
\subsection{Descripción del Negocio}

El sistema FoT tiene como propósito proveer a los agricultores (\textit{Farmers}) de una
plataforma de bajo costo para el monitoreo de variables ambientales y la gestión automatizada
del riego en granjas inteligentes. El actor principal del sistema es el \textbf{Farmer}: persona
encargada de operar y supervisar la finca, con conocimientos técnicos básicos o nulos. El sistema
opera en condiciones de restricción energética y conectividad limitada, propias del contexto
agrícola cubano, por lo que toda la infraestructura de procesamiento, almacenamiento y
visualización reside en la propia finca, sin dependencia de servicios externos ni de Internet.

% ─────────────────────────────────────────────────────────────────────────────
\subsection{Requisitos Funcionales}

Los requisitos funcionales describen las capacidades observables que el sistema debe ofrecer
al Farmer. Se organizan en dos grupos: el núcleo operativo del sistema y las funcionalidades
de gestión multi-parcela.

\begin{description}
	
	\item[RF-01 — Activar Riego Manualmente.]
	El sistema debe permitir al Farmer activar el sistema de riego (bomba de agua) de forma
	manual a través de la interfaz de usuario, independientemente de las lecturas de los
	sensores.
	
	\item[RF-02 — Detener Riego Manualmente.]
	El sistema debe permitir al Farmer desactivar el sistema de riego de forma manual en
	cualquier momento a través de la interfaz de usuario.
	
	\item[RF-03 — Activar/Desactivar Modo Riego Automático.]
	El sistema debe permitir al Farmer activar o desactivar un modo de funcionamiento
	automático, en el que el riego se gestione sin intervención manual directa.
	
	\item[RF-04 — Configurar Umbrales para Riego Automático.]
	Cuando el modo automático está activado, el sistema debe regar en función de umbrales
	configurables por el Farmer (p.\,ej.\ ``regar si la humedad del suelo es menor a
	\(X\%\)'').
	
	\item[RF-05 — Consultar Histórico de Variables.]
	El sistema debe permitir al Farmer consultar y visualizar en la interfaz de computadora
	el histórico almacenado de una o varias variables de medición en un rango de tiempo
	seleccionado.
	
	\item[RF-06 — Gestionar Parcelas.]
	El sistema debe permitir al Farmer \textbf{crear}, \textbf{modificar}, \textbf{eliminar}
	y \textbf{listar} las parcelas o zonas de cultivo que componen su explotación. Cada
	parcela tendrá un nombre único y, opcionalmente, una descripción, ubicación geográfica
	y tipo de cultivo.
	
	\item[RF-07 — Gestionar Sensores/Actuadores por Parcela.]
	El sistema debe permitir al Farmer \textbf{añadir}, \textbf{quitar}, \textbf{renombrar}
	y \textbf{reasignar} dispositivos (sensores y actuadores) a una parcela específica. Cada
	dispositivo tendrá un identificador único (p.\,ej.\ dirección MAC o ID lógico) y un tipo
	(p.\,ej.\ ``Sensor de Humedad'', ``Válvula de Riego'').
	
	\item[RF-08 — Visualizar Variables Agrupadas por Parcela.]
	El sistema debe mostrar las variables de medición organizadas por parcela. El Farmer
	podrá seleccionar una parcela y ver en tiempo real el estado de todos los sensores
	asociados a ella.
	
	\item[RF-09 — Control de Riego por Parcela.]
	Las acciones de activación/desactivación de riego, tanto manual como automática, deben
	poder realizarse de forma independiente para cada parcela.
	
	\item[RF-10 — Configurar Umbrales por Parcela.]
	La configuración de umbrales para el riego automático debe ser específica para cada
	parcela, dado que distintos cultivos o tipos de suelo presentan necesidades hídricas
	diferentes.
	
	\item[RF-11 — Consultar Histórico por Parcela y Variable.]
	El sistema debe permitir al Farmer consultar el histórico de cualquier variable, con
	posibilidad de filtrar por parcela, por sensor y por rango de fechas.
	
	\item[RF-12 — Detectar y Registrar Nuevos Dispositivos.]
	La estación base debe ser capaz de detectar automáticamente nuevos dispositivos
	(Arduinos con sensores) que se conecten a la red y guiar al Farmer en su asignación
	a una parcela.
	
	\item[RF-13 — Visualizar Mapa de la Finca \textnormal{(Opcional)}.]
	El sistema podrá ofrecer una vista de mapa, condicionada a la disponibilidad de
	coordenadas GPS, mostrando la ubicación de las parcelas y los dispositivos. Su
	implementación queda sujeta a disponibilidad de hardware GPS y se considera trabajo
	futuro.
	
\end{description}

% ─────────────────────────────────────────────────────────────────────────────
\subsection{Atributos de Calidad (Requisitos No Funcionales)}

Los atributos de calidad establecen las condiciones bajo las cuales el sistema debe ejecutar
sus funciones. Su definición responde a las restricciones del contexto cubano descritas en
el Capítulo~\ref{chap:1}.

\begin{description}
	
	\item[RNF-01 — Confiabilidad.]
	El sistema de riego automático debe activarse y desactivarse de forma fiable ante las
	condiciones de humedad definidas, evitando falsos positivos que provoquen anegamiento y
	falsos negativos que causen estrés hídrico. El margen de error en la lectura de sensores
	no debe superar el \(5\%\), valor compatible con la precisión especificada del sensor
	DHT22 (\(\pm0{,}5\,^\circ\)C; \(\pm2\%\) HR).
	
	\item[RNF-02 — Usabilidad.]
	La interfaz de usuario, tanto local (pantalla y botones) como remota, debe ser intuitiva
	para un agricultor con conocimientos técnicos básicos o nulos. El sistema debe poder
	operarse sin necesidad de consultar un manual extenso.
	
	\item[RNF-03 — Eficiencia Energética y Autonomía.]
	El sistema debe operar con alimentación inestable, soportando cortes eléctricos 
	mediante batería de respaldo. La autonomía mínima objetivo es de \textbf{24~horas}
	sin recarga, dimensionada para los ciclos de apagón programado del contexto cubano.
	El diseño debe ser compatible con fuentes alternativas de energía, como paneles
	solares de bajo costo.
	
	\item[RNF-04 — Mantenibilidad.]
	El software del sistema (código de Arduino y aplicación de escritorio) debe estar bien
	estructurado y documentado para facilitar correcciones, mejoras o adaptaciones a otros
	cultivos y sensores. Se exigirá código fuente comentado y uso de bibliotecas estándar
	de código abierto.
	
	\item[RNF-05 — Escalabilidad.]
	La arquitectura debe permitir la incorporación futura de nuevos sensores (p.\,ej.\ de
	salinidad o pH) o módulos de comunicación sin requerir un rediseño completo del sistema.
	
	\item[RNF-06 — Portabilidad.]
	La aplicación de monitoreo en computadora debe ejecutarse en equipos de bajos recursos,
	habituales en el contexto cubano, incluyendo hardware con más de diez años de
	antigüedad. Debe ser compatible con Windows y Linux, con consumo de RAM y CPU mínimos.
	
	\item[RNF-07 — Seguridad de Acceso.]
	El sistema debe garantizar que solo el Farmer autorizado pueda activar o detener el
	riego manualmente y modificar la configuración de umbrales. El mecanismo de
	autenticación será una contraseña local almacenada en la EEPROM del nodo o en un
	archivo de configuración en la estación base.
	
	\item[RNF-08 — Flexibilidad.]
	Los umbrales de riego y parámetros de configuración deben poder modificarse sin
	reiniciar el sistema completo; los cambios se persistirán en la EEPROM del nodo.
	Las modificaciones en la topología de parcelas o en la asignación de sensores pueden
	requerir el reinicio únicamente del nodo afectado.
	
	\item[RNF-09 — Tolerancia a Fallos.]
	La caída de un nodo sensor en una parcela no debe afectar el funcionamiento del resto
	del sistema ni la visualización de datos de las demás parcelas. El sistema debe registrar
	el fallo y notificarlo al Farmer en la interfaz.
	
\end{description}

% ─────────────────────────────────────────────────────────────────────────────
\subsection{Restricciones}

Las restricciones delimitan el espacio de solución válido y no son negociables, dado que
derivan directamente de las condiciones económicas, energéticas y tecnológicas del contexto
cubano.

\begin{description}
	
	\item[RES-01 — Plataforma hardware principal.]
	El sistema estará basado en la plataforma Arduino UNO por su bajo costo, disponibilidad
	en mercados alternativos y facilidad de programación.
	
	\item[RES-02 — Componentes de fácil reposición.]
	Todos los componentes deben ser de bajo costo y, en la medida de lo posible, de fácil
	adquisición en el mercado local cubano, reduciendo la dependencia de importaciones
	especializadas.
	
	\item[RES-03 — Conectividad limitada o nula.]
	El sistema debe funcionar correctamente en entornos con conectividad a Internet limitada
	o inexistente. El almacenamiento local mediante tarjeta SD es obligatorio e implementa
	el patrón \textit{store-and-forward} definido en el Capítulo~\ref{chap:1}.
	
	\item[RES-04 — Alimentación inestable.]
	El sistema debe ser capaz de operar con suministro eléctrico inestable o mediante
	baterías, considerando los apagones programados del contexto cubano.
	
	\item[RES-05 — Condiciones ambientales de campo.]
	Los componentes electrónicos deben soportar las condiciones ambientales del entorno
	agrícola: polvo, humedad elevada y temperaturas de hasta \(40\,^\circ\)C.
	
	\item[RES-06 — Software de código abierto.]
	El software para la estación base se desarrollará con lenguajes y herramientas de código
	abierto y multiplataforma (p.\,ej.\ Python), sin dependencias de licencias comerciales
	ni servicios de pago.
	
	\item[RES-07 — Nodo de campo por parcela.]
	Cada parcela dispondrá de su propio nodo de campo basado en Arduino UNO, con sus
	sensores y actuadores locales. Un único nodo puede gestionar múltiples sensores
	conectados a sus entradas analógicas y digitales.
	
	\item[RES-08 — Comunicación inalámbrica priorizada.]
	La comunicación entre los nodos de campo y la estación base será inalámbrica,
	priorizando \textbf{LoRa} para zonas sin infraestructura de red (alcance superior a
	15~km, consumo ultrabajo), \textbf{RF 433\,MHz} como alternativa de bajo costo cuando
	el alcance requerido sea menor, y \textbf{WiFi} únicamente cuando exista disponibilidad
	local. La dependencia de Internet es nula.
	
	\item[RES-09 — Independencia de servicios externos.]
	El sistema completo debe funcionar sin conexión a Internet ni servicios en la nube.
	Todo el procesamiento, almacenamiento y visualización ocurrirá en la infraestructura
	local del Farmer, diferenciando esta solución de plataformas como ThingSpeak
	analizadas en el Capítulo~\ref{chap:1}.
	
	\item[RES-10 — Resiliencia ante cortes eléctricos.]
	El sistema debe retomar operación automáticamente tras la restauración del suministro
	eléctrico, sin intervención del Farmer. Para ello se empleará el \textit{watchdog timer}
	del microcontrolador y el arranque automático de todos los servicios de la estación
	base.
	
\end{description}









% 
%  ############################################################
%  #                                                          #
%  #     RRRRR   OOO   U   U   N N N   DDDD      2222         #
%  #     R   R  O   O  U   U   N N N   D   D    2   2         #
%  #     RRRR   O   O  U   U   N N N   D   D      22          #
%  #     R R    O   O  U   U   N N N   D   D     2            #
%  #     R  R   OOO     UUU    N N N   DDDD     22222         #
%  #                                                          #
%  ############################################################
%







% ==============================================================
\section{Documentación de Diseño de Software}
\label{sec:disenio}
% ==============================================================

La documentación de diseño de software describe la arquitectura candidata
seleccionada para el sistema FoT, los componentes que la integran y los
patrones de diseño aplicados. Las decisiones adoptadas responden
directamente a las restricciones y atributos de calidad definidos en la
sección~\ref{sec:arq-def}: bajo costo, operación autónoma ante cortes
eléctricos, independencia de Internet y extensibilidad ante nuevos
sensores. A continuación se presentan la definición y descripción de la
arquitectura candidata, seguidas de los mecanismos de diseño identificados.


% --------------------------------------------------------------
\subsection{Definición de la Arquitectura Candidata}
\label{sec:arq-def}
% --------------------------------------------------------------

Para seleccionar la arquitectura candidata se evaluaron tres alternativas,
descritas en la tabla~\ref{tab:arq-candidatas}.

\begin{table}[htbp]
\centering
\caption{Evaluación de arquitecturas candidatas para el sistema FoT.}
\label{tab:arq-candidatas}
\small
\begin{tabular}{|p{3.4cm}|p{5.4cm}|p{4.6cm}|}
\hline
\textbf{Alternativa} & \textbf{Descripción} & \textbf{Veredicto} \\
\hline
Centralizada en la nube
    & Nodos envían todos los datos a un servidor remoto (ThingSpeak, AWS IoT
      u similar). Decisiones de riego tomadas en la nube.
    & \textbf{Rechazada.} Viola RES-03, RES-09 y RES-04: requiere
      conectividad continua a Internet, incompatible con el contexto cubano. \\
\hline
Peer-to-peer entre nodos
    & Los propios nodos Arduino intercambian decisiones entre sí sin un
      coordinador central.
    & \textbf{Rechazada.} El Arduino UNO carece de capacidad de cómputo
      para actuar como coordinador (RNF-06) y la arquitectura dificulta
      la gestión multi-parcela (RF-06--RF-11). \\
\hline
\textbf{Edge-first en tres capas}
    & Nodo de campo autónomo (lógica local), protocolo de mensajería
      ligero (MQTT) y estación base como coordinador y visualizador
      dentro de la red local de la finca.
    & \textbf{Seleccionada.} Satisface RES-03, RES-04, RES-08, RES-09
      y RNF-03 sin comprometer RF-05--RF-12. \\
\hline
\end{tabular}
\end{table}

La arquitectura \textit{edge-first} en tres capas combina tres estilos
arquitectónicos complementarios. El estilo en \textbf{capas} (\textit{Layered
Architecture}) estructura el sistema en tres niveles con dependencias
unidireccionales: Nodo de Campo $\to$ Comunicación $\to$ Estación Base. El
estilo \textbf{orientado a eventos} (\textit{Event-Driven Architecture})
gobierna el comportamiento dinámico: cada lectura de sensor es un evento que
puede desencadenar una decisión de riego o una actualización de la interfaz
sin necesidad de sondeo periódico (\textit{polling}). El estilo
\textbf{basado en componentes} (\textit{Component-Based Architecture})
organiza la capa de presentación de la estación base en módulos
independientes de interfaz gráfica (Python/Tkinter), gestión de datos y
lógica de negocio, facilitando el mantenimiento contemplado en RNF-04.


% --------------------------------------------------------------
\subsection{Descripción de la Arquitectura Candidata}
\label{sec:arq-desc}
% --------------------------------------------------------------

La figura~\ref{fig:arq-candidata} (representación esquemática en forma de
tabla) ilustra la arquitectura candidata seleccionada. El sistema se
estructura en tres capas con responsabilidades claramente delimitadas.

\begin{figure}[htbp]
\centering
\caption{Arquitectura candidata del sistema FoT (edge-first, tres capas).}
\label{fig:arq-candidata}
\small
\renewcommand{\arraystretch}{1.5}
\begin{tabular}{|c|p{11cm}|}
\hline
\textbf{Capa} & \textbf{Componentes y patrones} \\
\hline
\multirow{4}{*}{\parbox{2.5cm}{\centering Nodo de Campo\\Arduino UNO R3}}
 & Sensores físicos: DHT22 (temp./humedad), capacitivo (humedad suelo), RC522 RFID \\
 & \textbf{Adapter + Abstract Factory}: interfaz uniforme \texttt{ISensor}, \texttt{SensorFactory} \\
 & \textbf{State Machine}: estados \textit{Idle}, \textit{Monitoring}, \textit{Irrigating}, \textit{Fault} \\
 & \textbf{Memento}: configuración persistente en EEPROM (recuperación tras corte) \\
\hline
\multirow{2}{*}{\parbox{2.5cm}{\centering Comunicación}}
 & Medio: WiFi / LoRa / RF 433 MHz (según disponibilidad, RES-08) \\
 & \textbf{Broker MQTT} (Eclipse Mosquitto) en estación base, tópicos \texttt{fot/<parcela>/...} \\
\hline
\multirow{4}{*}{\parbox{2.5cm}{\centering Estación Base\\PC del Farmer\\Python 3 + SQLite}}
 & \textbf{Observer}: motor MQTT como sujeto, UI y motor de decisión como observadores \\
 & \textbf{Composite}: jerarquía \textit{Finca → Parcela → Dispositivo} \\
 & \textbf{Memento}: snapshot de configuración de parcelas \\
 & Base de datos SQLite, interfaz Tkinter, gestor de acceso (RNF-07) \\
\hline
\end{tabular}
\end{figure}

\noindent\textbf{Capa 1 — Nodo de Campo (Arduino UNO R3).} Constituye el
núcleo operativo del sistema. Adquiere lecturas de los sensores DHT22 y
capacitivo de humedad, las procesa localmente y activa el relay de la bomba
de agua según la lógica configurada. Opera de forma autónoma mediante una
máquina de estados finitos (patrón \textbf{State}) que contempla cuatro
estados: \textit{Idle} (modo manual o sin configuración), \textit{Monitoring}
(modo automático, lectura periódica de sensores), \textit{Irrigating} (bomba
activa) y \textit{Fault} (sensor fuera de rango o fallo de comunicación).
La autenticación física del Farmer se realiza a través del módulo RC522
RFID, eliminando la necesidad de teclado en campo (RNF-07). La configuración
de umbrales y el modo de operación se persisten en la EEPROM del
microcontrolador mediante el patrón \textbf{Memento}, lo que garantiza la
recuperación automática tras un corte eléctrico (RES-10).

Los sensores físicos se acceden a través de una capa de abstracción
implementada con el patrón \textbf{Adapter}, que expone una interfaz uniforme
\texttt{ISensor} independientemente del tipo de sensor subyacente. La
instanciación de estos adaptadores corresponde al patrón \textbf{Abstract
Factory}: la clase \texttt{SensorFactory} recibe el tipo de sensor como
parámetro y retorna el adaptador adecuado, lo que permite incorporar nuevos
sensores (p.ej., de salinidad o pH, según RNF-05) sin modificar el código
existente.

\noindent\textbf{Comunicación — MQTT sobre red local.} El protocolo MQTT,
descrito en la sección~1.3.3, actúa como capa de mensajería asíncrona entre
el nodo de campo y la estación base. El broker Eclipse Mosquitto se ejecuta
en el propio equipo del Farmer, eliminando toda dependencia de servicios
externos (RES-09). El nodo publica en los tópicos
\texttt{fot/<parcela>/sensores} y \texttt{fot/<parcela>/estado}, y se
suscribe a \texttt{fot/<parcela>/control} para recibir comandos. El uso de
QoS~1 garantiza la entrega al menos una vez, tolerando retransmisiones en
redes inestables. La prioridad de medio físico sigue el orden establecido en
RES-08: LoRa para extensiones rurales sin infraestructura, RF 433\,MHz como
alternativa de bajo costo, y WiFi cuando existe disponibilidad local.

\noindent\textbf{Capa 3 — Estación Base (Python~3 + SQLite).} La aplicación
de escritorio desarrollada en Python~3 con Tkinter recibe los eventos MQTT,
los procesa y los persiste en la base de datos SQLite. La notificación de
nuevos datos a la interfaz gráfica se implementa mediante el patrón
\textbf{Observer}: el motor de eventos MQTT actúa como sujeto (\textit{Subject})
y los componentes de visualización y el motor de decisión actúan como
observadores (\textit{Observers}). La jerarquía de entidades del dominio —finca,
parcelas y dispositivos— se modela con el patrón \textbf{Composite}, que
permite aplicar operaciones de lectura, filtrado y control de forma uniforme
a cualquier nivel de la jerarquía (RF-08, RF-09, RF-11). La configuración de
parcelas se gestiona mediante una segunda instancia del patrón
\textbf{Memento}, que permite deshacer cambios accidentales en la topología
de la finca. La aplicación está diseñada para ejecutarse en hardware con más
de diez años de antigüedad, tanto en Linux como en Windows (RNF-06), sin
dependencias externas que requieran licencias comerciales (RES-06).


% --------------------------------------------------------------
\subsection{Mecanismos de diseño identificados}
\label{sec:mecanismos}
% --------------------------------------------------------------

Los mecanismos de diseño establecen la correspondencia entre las necesidades
de análisis del sistema —expresadas como atributos de calidad y
restricciones— y las soluciones de diseño e implementación que las satisfacen.
La tabla~\ref{tab:mecanismos} resume los nueve mecanismos identificados para
el sistema FoT.

\begin{table}[htbp]
\centering
\caption{Mecanismos de diseño del sistema FoT.}
\label{tab:mecanismos}
\small
\begin{tabular}{|c|p{3.2cm}|p{3.6cm}|p{4.8cm}|}
\hline
\textbf{N\raisebox{0.5ex}{\textdegree}} & \textbf{Mecanismo de análisis} & \textbf{Mecanismo de diseño}
    & \textbf{Mecanismo de implementación} \\
\hline
1 & Persistencia de series temporales
    & Base de datos embebida con propiedades ACID
    & SQLite~3.x; módulo \texttt{sqlite3} de Python~3; sin servidor
      ni configuración de red (RNF-09, RES-06) \\
\hline
2 & Comunicación ligera entre nodo y estación
    & Pub/Sub asíncrono con broker local
    & Eclipse Mosquitto; biblioteca \texttt{PubSubClient} para Arduino;
      tópicos \texttt{fot/<parcela>/<función>} (RES-08, RES-09) \\
\hline
3 & Autenticación física sin teclado en campo
    & Identificación por RFID con almacenamiento local de credencial
    & Módulo RC522; biblioteca \texttt{MFRC522}; UID almacenado
      en EEPROM (RNF-07) \\
\hline
4 & Recuperación automática ante cortes eléctricos
    & Snapshot de configuración + arranque autónomo de servicios
    & Patrón \textbf{Memento}; EEPROM del ATmega328P; \textit{watchdog timer}
      integrado; \texttt{systemd} (Linux) o Tarea Programada (Windows)
      para reinicio de servicios (RES-04, RES-10) \\
\hline
5 & Extensibilidad sin modificación del núcleo
    & Abstracción de dispositivos mediante interfaz uniforme
    & Patrones \textbf{Adapter} + \textbf{Abstract Factory};
      interfaz \texttt{ISensor}; clase \texttt{SensorFactory}
      configurable por archivo de parámetros (RNF-05) \\
\hline
6 & Gestión jerárquica de entidades del dominio
    & Composición recursiva de entidades
    & Patrón \textbf{Composite}; clases \texttt{Finca},
      \texttt{Parcela}, \texttt{Dispositivo} con interfaz
      \texttt{IComponente} (RF-06, RF-07, RF-08) \\
\hline
7 & Notificación de cambios en tiempo real a la UI
    & Desacoplamiento productor-consumidor de eventos
    & Patrón \textbf{Observer}; motor de eventos MQTT como
      \textit{Subject}; módulos de UI y decisión como
      \textit{Observers} (RNF-01, RF-08) \\
\hline
8 & Control robusto de la lógica de irrigación
    & Máquina de estados finitos explícita
    & Patrón \textbf{State}; estados \textit{Idle},
      \textit{Monitoring}, \textit{Irrigating}, \textit{Fault};
      transiciones validadas en el nodo de campo (RF-03, RNF-01) \\
\hline
9 & Autonomía energética del nodo de campo
    & Modo de bajo consumo con reactivación por interrupción
    & Modos \textit{sleep} del ATmega328P;
      \textit{watchdog timer} como temporizador de ciclo;
      consumo en reposo inferior a 10\,mA (RNF-03, RES-04) \\
\hline
\end{tabular}
\end{table}


% ==============================================================
\section{Conclusiones parciales}
\label{sec:cap2-concl}
% ==============================================================

Una vez concluida la documentación de ingeniería de requisitos y de diseño
de software del sistema FoT, se arriban a las siguientes conclusiones:

\begin{enumerate}
\setlength\itemsep{0.3em}

\item La documentación de requisitos funcionales, atributos de calidad y
restricciones evidencia que el contexto cubano no constituye únicamente un
conjunto de limitaciones a mitigar, sino el criterio de diseño principal: la
independencia de Internet (RES-09), la tolerancia a cortes eléctricos
(RES-04, RES-10) y el empleo exclusivo de componentes de bajo costo y código
abierto (RES-01, RES-06) definen un espacio de solución bien delimitado que
descarta las alternativas centralizadas en la nube y peer-to-peer evaluadas
en la sección~\ref{sec:arq-def}.

\item La arquitectura edge-first en tres capas, con los estilos en capas,
orientado a eventos y basado en componentes, satisface simultáneamente los
requisitos de operación autónoma local (RNF-03, RES-04), comunicación
inalámbrica sin dependencia de Internet (RES-08, RES-09) y mantenibilidad
en hardware de bajos recursos (RNF-06), aspectos que ninguna de las
soluciones similares revisadas en el capítulo~1 aborda de forma integral.

\item Los seis patrones GoF identificados —\textbf{Adapter}, \textbf{Abstract
Factory}, \textbf{Observer}, \textbf{Composite}, \textbf{State} y
\textbf{Memento}— responden cada uno a un problema concreto del dominio:
la extensibilidad sin modificación del núcleo (Adapter + Abstract Factory),
el desacoplamiento entre adquisición de datos y visualización (Observer), la
gestión uniforme de la jerarquía finca-parcela-dispositivo (Composite), la
robustez de la lógica de irrigación autónoma (State) y la recuperación de
configuración tras cortes eléctricos (Memento). Su aplicación conjunta
garantiza que los atributos de calidad RNF-01, RNF-05, RNF-08 y RNF-09 se
encuentren respaldados por soluciones de diseño documentadas y verificables.

\item Los nueve mecanismos de diseño establecidos en la
sección~\ref{sec:mecanismos} vinculan explícitamente cada necesidad de
análisis con su solución de implementación concreta, proporcionando la
trazabilidad necesaria para que la validación del capítulo~3 pueda
contrastar el comportamiento observado del prototipo con los compromisos de
diseño aquí adquiridos.

\end{enumerate}

\pagebreak